% Diagram of the top of the proof (Introduction): from the final theorems down to Theorems A, B, C.
\begin{figure}[tbp]
\centering
\begin{tikzpicture}[x=1cm,y=1cm]
% row 0: the statements on the critical line
\node[pfmain] (line) at (0,0) {Thms.~\ref{intro:thm:line}, \ref{intro:thm:rh}\\zeros on the critical line};
% row 1
\node[pfmain] (grh) at (0,-1.7) {Thm.~\ref{main:thm:grh}\\(=~Thm.~\ref{intro:thm:grh})\\half-plane theorem};
\node[pfnode] (fe) at (6,-1.7) {Prop.~\ref{zeros:prop:line}\\functional equation};
% row 2
\node[pfnode] (crit) at (-3.75,-3.4) {Thm.~\ref{zeros:thm:criterion}\\one-sided criterion};
\node[pfnode] (one) at (3.75,-3.4) {Thm.~\ref{main:thm:onesided}\\one-sided bound};
% row 3
\node[pffact, text width=1.9cm] (landau) at (-6.25,-5.1) {Fact~\ref{zeros:fact:landau}\\Landau's lemma};
\node[pffact, text width=1.9cm] (prime) at (-3.75,-5.1) {Fact~\ref{zeros:fact:prime}\\choice of the prime};
\node[pfnode, text width=1.9cm] (poles) at (-1.25,-5.1) {Prop.~\ref{zeros:prop:poles}\\poles};
\node[pfnode, text width=1.9cm] (atx) at (1.25,-5.1) {Prop.~\ref{main:prop:atx}\\remainder at the first scale};
\node[pfnode, text width=1.9cm] (exist) at (3.75,-5.1) {Thm.~\ref{setup:thm:existence}\\existence of the data};
\node[pfnode, text width=1.9cm] (pnt) at (6.25,-5.1) {Thm.~\ref{zeros:thm:pnt}\\prime number theorem};
% row 4
\node[pfnode] (rem) at (-1.75,-6.8) {Prop.~\ref{setup:prop:rem}\\total remainder};
\node[pfmain] (star) at (1.25,-6.8) {Thm.~\ref{main:thm:star}\\main inequality};
% row 5
\node[pfmain] (A) at (-1.75,-8.5) {Thm.~\ref{main:thm:A}\\Theorem A: upper bound};
\node[pfmain] (B) at (1.25,-8.5) {Thm.~\ref{main:thm:B}\\Theorem B: spanning};
\node[pfmain] (C) at (4.25,-8.5) {Thm.~\ref{main:thm:C}\\Theorem C: lower bound};
% row 6: the parts
\node[pfpart] (P2) at (-1.75,-10.2) {Part~2\\seesaw identity,\\divisibility, size\\Fig.~\ref{fig:part2}};
\node[pfpart] (P3) at (1.25,-10.2) {Part~3\\second moment\\Fig.~\ref{fig:part3}};
\node[pfpart] (P4) at (4.25,-10.2) {Part~4\\chain of lifts,\\twists\\Figs.~\ref{fig:part4a}, \ref{fig:part4b}};
% arrows: premise -- conclusion
\draw[pfedge] (grh) -- (line);
\draw[pfedge] (fe) -- (line);
\draw[pfedge] (crit) -- (grh);
\draw[pfedge] (one) -- (grh);
\draw[pfedge] (landau) -- (crit);
\draw[pfedge] (prime) -- (crit);
\draw[pfedge] (poles) -- (crit);
\draw[pfedge] (atx) -- (one);
\draw[pfedge] (exist) -- (one);
\draw[pfedge] (pnt) -- (one);
\draw[pfedge] (rem) -- (atx);
\draw[pfedge] (star) -- (atx);
\draw[pfedge] (A) -- (star);
\draw[pfedge] (B) -- (star);
\draw[pfedge] (C) -- (star);
\draw[pfedge] (P2) -- (A);
\draw[pfedge] (P3) -- (B);
\draw[pfedge] (P4) -- (C);
\end{tikzpicture}
\caption{The top of the proof: from the theorems of the introduction down to Theorems~A, B
and~C (Sections~\ref{sec:intro}, \ref{sec:zeros}, \ref{sec:setup} and~\ref{sec:main}). Each box
is a result, with its number and a short name; a box with several numbers stands for a group of
results. An arrow from one box to another means that the proof of the second, as written in this
paper, uses the first. Bold boxes are the final theorems,
the main inequality and Theorems~A, B and~C. Dashed grey boxes are background results admitted
without proof. Double boxes stand for the parts that prove Theorems~A, B and~C; each has its own
diagram.}
\label{fig:overview}
\end{figure}
